\documentclass[12pt,a4paper]{article}
\usepackage{iftex}
\ifPDFTeX
  \usepackage[T2A]{fontenc}
  \usepackage[utf8]{inputenc}
  \usepackage[english,russian]{babel}
\else
  \usepackage{fontspec}
  \setmainfont{DejaVu Serif}
  \setmonofont{DejaVu Sans Mono}
  \renewcommand{\contentsname}{Содержание}
  \renewcommand{\figurename}{Рис.}
  \renewcommand{\tablename}{Таблица}
\fi

\usepackage[a4paper,margin=25mm]{geometry}
\usepackage{amsmath,amssymb,amsthm}
\usepackage{graphicx}
\usepackage{booktabs}
\usepackage{enumitem}
\usepackage[hidelinks]{hyperref}

\newcommand{\Exp}{\mathsf{M}}
\newcommand{\Var}{\mathsf{D}}
\newcommand{\PP}{\mathsf{P}}
\newcommand{\dens}{\mathrm{p}}
\newcommand{\Normal}{\mathcal{N}}
\newcommand{\R}{\mathbb{R}}
\newcommand{\xs}[1]{x^{(#1)}}
\newcommand{\ys}[1]{y^{(#1)}}
\newcommand{\norm}[1]{\left\lVert #1 \right\rVert}
\newcommand{\abs}[1]{\left\lvert #1 \right\rvert}
\newcommand{\dd}{\,\mathrm{d}}
\newcommand{\est}[1]{\tilde{#1}}
\newcommand{\zad}[1]{\medskip\noindent\textbf{Задание #1.}\quad}
\newcommand{\prov}{\smallskip\noindent\textit{Проверка:}\ }

\theoremstyle{definition}
\newtheorem{definition}{Определение}
\newtheorem{statement}{Утверждение}
\theoremstyle{plain}
\newtheorem{lemma}{Лемма}

% Список с русскими буквенными метками
\newenvironment{ruslist}
  {\begin{itemize}[label={},leftmargin=2.2em,itemsep=2pt,topsep=3pt]}
  {\end{itemize}}


\begin{document}

\begin{center}
  {\large\bfseries Задача 1. Элементы теории вероятностей}\\[1mm]
  Нейронные сети и машинное обучение\\[1mm]
  Задания 1--9
\end{center}

\section{Что нужно знать}

\subsection{Обозначения}

\begin{center}
\begin{tabular}{ll}
\toprule
$\xi,\ \eta,\ \zeta$ & случайные величины \\
$\PP(A)$ & вероятность события $A$ \\
$F_\xi(x) = \PP(\xi \le x)$ & функция распределения \\
$\dens_\xi(x)$ & плотность распределения \\
$\Exp\xi$, $\Var\xi$ & математическое ожидание и дисперсия \\
$\mathrm{cov}(\xi,\eta)$ & ковариация \\
$\xi \sim \mathrm{U}[a,b]$ & равномерное распределение на отрезке \\
$\xi \sim \Normal(\mu,\sigma^2)$ & нормальное распределение \\
$\varphi$, $\Phi$ & плотность и функция распределения $\Normal(0,1)$ \\
$\Sigma = (\sigma_{ij})$ & ковариационная матрица \\
п.\,н. & почти наверное, то есть с вероятностью 1 \\
\bottomrule
\end{tabular}
\end{center}

\subsection{Функция распределения и плотность}

\emph{Функция распределения} $F_\xi(x) = \PP(\xi \le x)$ отвечает на вопрос:
какова вероятность, что случайная величина не превысит $x$. Она не убывает,
непрерывна справа, $F_\xi(-\infty) = 0$, $F_\xi(+\infty) = 1$. Функция
распределения задаёт распределение однозначно: если у двух случайных величин
совпали функции распределения, то совпали и распределения. Поэтому стандартный
способ доказать «$\eta$ имеет распределение $F$» -- посчитать $F_\eta$ и
убедиться, что получилось $F$.

\emph{Плотность} $\dens_\xi$ -- это производная функции распределения:
\[
  F_\xi(x) = \int_{-\infty}^{x}\dens_\xi(u)\dd u,
  \qquad \dens_\xi(x) = F_\xi'(x).
\]
Плотность неотрицательна и интегрируется в единицу. Чем больше $\dens_\xi(x)$,
тем чаще значения $\xi$ попадают в окрестность точки $x$.

Для $\xi \sim \mathrm{U}[0,1]$ плотность равна $1$ на $[0,1]$ и нулю вне
отрезка, а $F_\xi(u) = u$ при $u \in [0,1]$. Это свойство используется дальше
постоянно: \emph{для равномерной на $[0,1]$ величины вероятность попасть ниже
уровня $u$ равна самому $u$}.

Для $\Normal(\mu,\sigma^2)$
\[
  \dens(x) = \frac{1}{\sigma\sqrt{2\pi}}\,e^{-(x-\mu)^2/2\sigma^2},
  \qquad
  \varphi(u) = \frac{1}{\sqrt{2\pi}}e^{-u^2/2},
  \qquad
  \Phi(u) = \int_{-\infty}^{u}\varphi .
\]

\subsection{Математическое ожидание и дисперсия}

\[
  \Exp\xi = \int_{\R} x\,\dens_\xi(x)\dd x,
  \qquad
  \Var\xi = \Exp(\xi - \Exp\xi)^2 = \Exp\xi^2 - (\Exp\xi)^2 .
\]
Математическое ожидание -- это среднее значение, дисперсия -- мера разброса
вокруг него. Дальше нужны два свойства.

\emph{Линейность математического ожидания:}
$\Exp(a\xi + b\eta + c) = a\Exp\xi + b\Exp\eta + c$. Никакой независимости здесь
не требуется, это верно всегда.

\emph{Дисперсия суммы:}
$\Var(\xi + \eta) = \Var\xi + \Var\eta + 2\,\mathrm{cov}(\xi,\eta)$.
Слагаемое с ковариацией исчезает, если величины независимы, и тогда дисперсии
просто складываются.

\subsection{Независимость, ковариация, ковариационная матрица}

Величины $\xi$ и $\eta$ \emph{независимы}, если знание значения одной ничего не
говорит о другой; формально $\PP(\xi \le x,\ \eta \le y) = \PP(\xi \le x)\PP(\eta \le y)$
при всех $x, y$. Основное следствие, которым будем пользоваться:
\begin{equation}\label{eq:indep}
  \xi,\ \eta \text{ независимы} \;\Longrightarrow\; \Exp(\xi\eta) = \Exp\xi\cdot\Exp\eta .
\end{equation}

\emph{Ковариация} $\mathrm{cov}(\xi,\eta) = \Exp\bigl((\xi - \Exp\xi)(\eta - \Exp\eta)\bigr)$
измеряет, насколько согласованно величины отклоняются от своих средних. Для
случайного вектора $\xi = (\xi_1,\dots,\xi_m)$ все попарные ковариации собирают
в \emph{ковариационную матрицу}
\[
  \Sigma = (\sigma_{ij}), \qquad
  \sigma_{ij} = \Exp\bigl((\xi_i - \mu_i)(\xi_j - \mu_j)\bigr),
  \qquad \mu_i = \Exp\xi_i .
\]
На диагонали стоят дисперсии: $\sigma_{ii} = \Var\xi_i$. Матрица симметрична,
поскольку $\sigma_{ij} = \sigma_{ji}$.

\subsection{Многомерное нормальное распределение}

Для вектора $\xi = (\xi_1,\dots,\xi_m)$ нормальная плотность имеет вид
\begin{equation}\label{eq:mvn}
  \dens(x) = \frac{1}{(2\pi)^{m/2}\det(\Sigma)^{1/2}}
  \exp\Bigl(-\tfrac12 (x-\mu)^{\top}\Sigma^{-1}(x-\mu)\Bigr),
  \qquad \xi \sim \Normal(\mu, \Sigma),
\end{equation}
где $\mu$ -- вектор средних, $\Sigma$ -- ковариационная матрица. Формула
работает только при $\det\Sigma \ne 0$: иначе обратной матрицы нет и плотности
не существует.

\subsection{Мода и медиана}

\emph{Мода} -- точка максимума плотности, наиболее вероятное значение.
\emph{Медиана} -- такое $m$, что $\PP(\xi \le m) \ge 1/2$ и
$\PP(\xi \ge m) \ge 1/2$, то есть значение, делящее распределение пополам.
Медиана определяется только через сравнения «больше -- меньше», а
математическое ожидание -- через расстояния между значениями. Из-за этого они
по-разному ведут себя при преобразованиях величины, что и разбирается
в заданиях 4 и 5.

\subsection{Неравенство Йенсена}

Если функция $h$ выпукла (график лежит не выше хорды), то
\begin{equation}\label{eq:jensen}
  \Exp h(\xi) \ge h(\Exp\xi),
\end{equation}
причём при строго выпуклой $h$ и невырожденной $\xi$ неравенство строгое. Для
вогнутой $h$ знак противоположный. Равенство выполняется при всех
распределениях тогда и только тогда, когда $h$ аффинна, то есть
$h(x) = ax + b$. Смысл: среднее от функции и функция от среднего -- разные
вещи, если функция не прямая.

\subsection{Замена переменной в плотности}

Пусть $g$ монотонна и непрерывно дифференцируема, и величины связаны равенством
$\xi = g(\zeta)$. Тогда плотности связаны так:
\begin{equation}\label{eq:jac}
  \dens_\zeta(t) = \dens_\xi\bigl(g(t)\bigr)\cdot\abs{g'(t)} .
\end{equation}
Множитель $\abs{g'(t)}$ появляется потому, что замена растягивает или сжимает
ось: отрезок $\dd t$ переходит в отрезок $\abs{g'(t)}\dd t$, и чтобы полная
вероятность осталась равной единице, плотность приходится домножить. Именно этот
множитель ломает наивные ожидания в задании 4.

\subsection{Условная вероятность и полная группа событий}

\emph{Условная вероятность} $\PP(A \mid B) = \PP(A \cap B)/\PP(B)$ при
$\PP(B) > 0$ -- вероятность $A$ при известном $B$.

События $W_1,\dots,W_n$ образуют \emph{полную группу}, если они попарно
несовместны ($W_i \cap W_j = \varnothing$ при $i \ne j$) и вместе покрывают всё
пространство. Тогда любое событие $X$ разбивается на непересекающиеся куски
$X \cap W_i$, откуда \emph{формула полной вероятности}
\begin{equation}\label{eq:total}
  \PP(X) = \sum_{i=1}^{n}\PP(X \cap W_i) = \sum_{i=1}^{n}\PP(X \mid W_i)\PP(W_i).
\end{equation}
Вероятности $\PP(W_i)$ называют \emph{априорными} (до наблюдения $X$), а
$\PP(W_i \mid X)$ -- \emph{апостериорными} (после).

\section{Задания}

\zad{1}
\dano Случайная величина $\xi \sim \mathrm{U}[0,1]$ и произвольная функция
распределения $F$.

\treb Доказать, что случайная величина $\eta = F^{-1}(\xi)$ имеет функцию
распределения $F$.

\zachem Генератор случайных чисел в компьютере выдаёт равномерные числа на
$[0,1]$. Это утверждение говорит, как из них получить величину с любым нужным
законом: достаточно применить обратную функцию распределения.

\resh Разберём сначала простой случай: $F$ непрерывна и строго возрастает. Тогда
$F^{-1}$ -- обычная обратная функция, и она тоже строго возрастает. Строго
возрастающая функция сохраняет неравенства, поэтому
\[
  F^{-1}(\xi) \le x \iff \xi \le F(x),
\]
то есть события $\{\eta \le x\}$ и $\{\xi \le F(x)\}$ -- это одно и то же
событие. Значит
\[
  F_\eta(x) = \PP(\eta \le x) = \PP\bigl(\xi \le F(x)\bigr) = F(x),
\]
где последний переход -- это свойство равномерного распределения
($\PP(\xi \le u) = u$ при $u \in [0,1]$, а $F(x)$ как раз лежит в $[0,1]$).
Функции распределения совпали, значит совпали распределения. $\blacksquare$

\smallskip
\emph{Общий случай.} Если $F$ имеет скачки или участки постоянства, обычной
обратной функции нет, и вместо неё берут
\[
  F^{-1}(u) = \inf\{x \in \R : F(x) \ge u\}
\]
(наименьшее $x$, на котором $F$ дотягивает до уровня $u$). Всё доказательство
проходит без изменений, если проверить ту же равносильность.

\begin{lemma}
Для неубывающей непрерывной справа $F$ выполняется
$F^{-1}(u) \le x \iff u \le F(x)$.
\end{lemma}
\begin{proof}
Обозначим $A_u = \{y : F(y) \ge u\}$, тогда $F^{-1}(u) = \inf A_u$.

Пусть $u \le F(x)$. Тогда $x \in A_u$, а инфимум множества не больше любого его
элемента, значит $F^{-1}(u) \le x$.

Пусть $F^{-1}(u) \le x$. Множество $A_u$ -- луч, уходящий вправо: если
$y \in A_u$ и $y' > y$, то $F(y') \ge F(y) \ge u$, то есть $y' \in A_u$. Его
левый конец $t = F^{-1}(u)$ сам принадлежит $A_u$, потому что $F$ непрерывна
справа: $F(t) = \lim_{y \downarrow t}F(y) \ge u$. Из $t \le x$ и монотонности
$F$ получаем $F(x) \ge F(t) \ge u$.
\end{proof}

Обратим внимание: использовались только неубывание и непрерывность справа,
которые есть у любой функции распределения. Поэтому способ работает и для
дискретных распределений, где $F^{-1}$ реализуется поиском по таблице
накопленных вероятностей.

\zad{2}
\dano Число $\mu$, число $\sigma > 0$ и случайная величина $\xi \sim \Normal(0,1)$.

\treb Доказать, что $\eta = \mu + \sigma\xi$ распределена по закону
$\Normal(\mu,\sigma^2)$.

\zachem Вместе с заданием 1 это даёт полный рецепт: получить стандартную
нормальную величину, а затем сдвигом и растяжением сделать из неё нормальную
с любыми параметрами.

\resh Считаем функцию распределения $\eta$, раскрывая неравенство под знаком
вероятности:
\[
  F_\eta(x) = \PP(\mu + \sigma\xi \le x)
            = \PP\Bigl(\xi \le \frac{x-\mu}{\sigma}\Bigr)
            = \Phi\Bigl(\frac{x-\mu}{\sigma}\Bigr).
\]
Деление на $\sigma > 0$ знак неравенства не меняет. Дифференцируем по $x$;
производная сложной функции даёт множитель $1/\sigma$:
\[
  \dens_\eta(x) = \frac{1}{\sigma}\,\varphi\Bigl(\frac{x-\mu}{\sigma}\Bigr)
  = \frac{1}{\sigma\sqrt{2\pi}}\exp\Bigl(-\frac{(x-\mu)^2}{2\sigma^2}\Bigr).
\]
Это в точности плотность $\Normal(\mu,\sigma^2)$. $\blacksquare$

Параметры можно проверить и напрямую, по линейности:
$\Exp\eta = \mu + \sigma\Exp\xi = \mu$ и $\Var\eta = \sigma^2\Var\xi = \sigma^2$.

Условие $\sigma > 0$ существенно только для удобства: при $\sigma < 0$
получается та же плотность с $\abs{\sigma}$ в знаменателе, так как $\varphi$
чётна, а при $\sigma = 0$ величина $\eta$ постоянна и плотности у неё нет.

\zad{3}
\dano Независимые случайные величины $\xi_1,\dots,\xi_n$, каждая распределена
равномерно на $[0,1]$, и их сумма $S_n = \sum_{i=1}^{n}\xi_i$.

\treb Доказать, что $\Exp S_n = n/2$ и $\Var S_n = n/12$.

\zachem По центральной предельной теореме сумма многих независимых слагаемых
близка к нормальной. Зная $\Exp S_n$ и $\Var S_n$, сумму можно
отнормировать и получить приближение к $\Normal(0,1)$ -- это второй способ
генерировать нормальную величину.

\resh Сначала одно слагаемое. Плотность равна единице на $[0,1]$, поэтому
интегрирование идёт по этому отрезку:
\[
  \Exp\xi = \int_0^1 x\dd x = \frac12,
  \qquad
  \Exp\xi^2 = \int_0^1 x^2\dd x = \frac13,
  \qquad
  \Var\xi = \Exp\xi^2 - (\Exp\xi)^2 = \frac13 - \frac14 = \frac{1}{12}.
\]
Теперь сумма. По линейности математического ожидания
\[
  \Exp S_n = \sum_{i=1}^{n}\Exp\xi_i = n\cdot\frac12 = \frac{n}{2}
\]
(независимость здесь не нужна). Для дисперсии
\[
  \Var S_n = \sum_{i=1}^{n}\Var\xi_i + 2\sum_{i<j}\mathrm{cov}(\xi_i,\xi_j)
           = n\cdot\frac{1}{12} = \frac{n}{12},
\]
где все ковариации равны нулю по независимости (это доказано в задании 7).
$\blacksquare$

\smallskip
Отсюда получается нормировка: величина
\[
  \eta = \frac{S_n - n/2}{\sqrt{n/12}}
\]
имеет $\Exp\eta = 0$ и $\Var\eta = 1$, а по центральной предельной теореме её
распределение с ростом $n$ приближается к $\Normal(0,1)$. При $n = 12$
знаменатель равен единице, и формула упрощается до $\eta = S_{12} - 6$.

Приближение остаётся приближением: $S_{12}$ принимает значения только из
отрезка $[0,12]$, поэтому у $\eta$ нет хвостов за пределами $[-6,6]$, тогда как
у настоящей нормальной величины они есть. Для моделирования шума этого хватает,
для оценки вероятностей редких событий -- нет.

\zad{4}
\dano Случайная величина $\xi$ с плотностью $\dens_\xi$. Делается замена
координат $x = g(t)$, то есть вводится новая величина $\zeta$, связанная с
исходной равенством $\xi = g(\zeta)$. Её плотность вычисляется по правилу
\eqref{eq:jac}. Функция $g$ монотонна и непрерывно дифференцируема.

\treb Выяснить, всегда ли (а) $\Exp\xi = \Exp\zeta$ и (б) точка максимума
плотности переходит в точку максимума, то есть $x_0 = g(t_0)$, где $x_0$ и $t_0$
-- моды $\xi$ и $\zeta$.

\zachem Замена переменной -- это смена единиц измерения или шкалы признака,
например переход к логарифму. Вопрос в том, что при этом сохраняется, а что нет.

\resh Ответ на оба вопроса: нет, не всегда.

\emph{(а) Математическое ожидание.} По формуле замены переменной под интегралом
\[
  \Exp\xi = \int \dens_\xi(x)\,x\dd x
          = \int \dens_\zeta(t)\,g(t)\dd t
          = \Exp g(\zeta).
\]
То есть равенство $\Exp\xi = \Exp\zeta$ означало бы $\Exp g(\zeta) = \Exp\zeta$,
чего для произвольной $g$ нет. Более того, не выполняется даже более слабое
$\Exp\xi = g(\Exp\zeta)$: по неравенству Йенсена \eqref{eq:jensen} для выпуклой
$g$ имеем $\Exp g(\zeta) \ge g(\Exp\zeta)$, причём строго, если $g$ строго
выпукла. Равенство при всех распределениях возможно только для аффинной $g$.

\emph{(б) Мода.} Прологарифмируем \eqref{eq:jac}:
\[
  \ln\dens_\zeta(t) = \ln\dens_\xi\bigl(g(t)\bigr) + \ln\abs{g'(t)} .
\]
В точке максимума $t_0$ производная равна нулю:
\[
  0 = \frac{\dens_\xi'\bigl(g(t_0)\bigr)}{\dens_\xi\bigl(g(t_0)\bigr)}\,g'(t_0)
      + \frac{g''(t_0)}{g'(t_0)} .
\]
Чтобы точка $x_0 = g(t_0)$ была максимумом $\dens_\xi$, нужно
$\dens_\xi'(x_0) = 0$, а тогда из равенства выше следует $g''(t_0) = 0$. Для
нелинейной $g$ это в общем случае неверно: слагаемое $\ln\abs{g'}$ сдвигает
максимум.

\emph{Контрпример сразу для обоих пунктов.} Возьмём $\xi \sim \mathrm{Exp}(1)$,
то есть $\dens_\xi(x) = e^{-x}$ при $x \ge 0$, и $g(t) = t^3$, то есть
$\zeta = \xi^{1/3}$. По \eqref{eq:jac}
\[
  \dens_\zeta(t) = e^{-t^3}\cdot 3t^2, \qquad t \ge 0 .
\]
Считаем средние:
\[
  \Exp\xi = 1, \qquad
  \Exp\zeta = \int_0^{\infty} t\cdot 3t^2 e^{-t^3}\dd t
            = \Gamma\Bigl(\frac43\Bigr) \approx 0.893 \ne 1,
\]
и заодно $g(\Exp\zeta) = \Gamma(4/3)^3 \approx 0.712 \ne 1$, что и есть строгое
неравенство Йенсена для выпуклой $g(t) = t^3$.

Считаем моды. Плотность $e^{-x}$ убывает, её максимум в $x_0 = 0$. Для
$\dens_\zeta$ удобно максимизировать логарифм $2\ln t - t^3$:
\[
  \frac{2}{t} - 3t^2 = 0 \;\Longrightarrow\;
  t_0 = \Bigl(\frac23\Bigr)^{1/3} \approx 0.874,
  \qquad
  g(t_0) = t_0^3 = \frac23 \approx 0.667 \ne 0 = x_0 .
\]
Оба равенства нарушены. $\blacksquare$

\smallskip
Вывод: ни среднее, ни положение максимума плотности не сохраняются при
нелинейной замене координат, виноват якобиан $\abs{g'}$. Оба свойства
восстанавливаются только для аффинной $g$.

\zad{5}
\dano Случайная величина $\xi$ и функция $z$; новая величина $\zeta = z(\xi)$.

\treb Выяснить (а) всегда ли $\Exp\zeta = z(\Exp\xi)$ и (б) как связаны медианы
$\xi$ и $\zeta$.

\resh
\emph{(а) Среднее не переносится.} Контрпример: $\xi \sim \mathrm{U}[0,1]$,
$z(x) = x^2$. Тогда
\[
  \Exp\zeta = \int_0^1 x^2\dd x = \frac13,
  \qquad
  z(\Exp\xi) = \Bigl(\frac12\Bigr)^2 = \frac14,
\]
и $1/3 \ne 1/4$. В общем виде это неравенство Йенсена \eqref{eq:jensen}: для
строго выпуклой $z$ всегда $\Exp z(\xi) > z(\Exp\xi)$, для строго вогнутой знак
обратный. Равенство при всех распределениях выполняется тогда и только тогда,
когда $z$ аффинна: при $z(x) = ax + b$ оно следует из линейности
математического ожидания.

\emph{(б) Медиана переносится, если $z$ монотонна.} Пусть $z$ непрерывна и
строго возрастает, а $m$ -- медиана $\xi$. Строго возрастающая функция сохраняет
неравенства, поэтому события $\{\xi \le m\}$ и $\{z(\xi) \le z(m)\}$ совпадают.
Значит
\[
  \PP\bigl(\zeta \le z(m)\bigr) = \PP(\xi \le m) \ge \frac12,
  \qquad
  \PP\bigl(\zeta \ge z(m)\bigr) = \PP(\xi \ge m) \ge \frac12,
\]
то есть $z(m)$ удовлетворяет определению медианы для $\zeta$. Для строго
убывающей $z$ рассуждение то же, только неравенства меняются местами, и вывод
сохраняется.

Причина различия между пунктами (а) и (б) уже отмечалась: медиана определяется
только порядком на прямой, который монотонная функция сохраняет, а среднее
зависит от расстояний между значениями, которые нелинейная $z$ искажает.

\emph{Немонотонная $z$ медиану не переносит.} При $\xi \sim \mathrm{U}[0,1]$ и
$z(x) = (x - 1/2)^2$ медиана $\xi$ равна $1/2$, и $z(1/2) = 0$. Но
\[
  \PP(\zeta \le m) = \PP\Bigl(\abs{\xi - \tfrac12} \le \sqrt m\Bigr) = 2\sqrt m,
  \qquad
  2\sqrt m = \frac12 \;\Longrightarrow\; m = \frac{1}{16} \ne 0 .
\]

\zad{6}
\dano Вектор $\xi = (\xi_1,\dots,\xi_m)$ имеет многомерное нормальное
распределение с плотностью \eqref{eq:mvn}, все его компоненты взаимно
независимы.

\treb Показать, что плотность \eqref{eq:mvn} равна произведению одномерных
нормальных плотностей компонент.

\zachem Это ровно то предположение, на котором строится наивный байесовский
классификатор: при независимых признаках многомерная модель распадается на
$m$ одномерных, и число её параметров падает с $m + m(m+1)/2$ до $2m$.

\resh По заданию 7 из независимости следует $\sigma_{ij} = 0$ при $i \ne j$,
поэтому ковариационная матрица диагональна:
\[
  \Sigma = \mathrm{diag}(\sigma_1^2,\dots,\sigma_m^2),
  \qquad \sigma_i^2 = \Var\xi_i > 0 .
\]
Для диагональной матрицы определитель равен произведению диагональных элементов,
а обратная получается покомпонентным делением:
\[
  \det\Sigma = \prod_{i=1}^{m}\sigma_i^2,
  \qquad
  \Sigma^{-1} = \mathrm{diag}\bigl(\sigma_1^{-2},\dots,\sigma_m^{-2}\bigr).
\]
Тогда квадратичная форма в показателе экспоненты превращается в сумму:
\[
  (x-\mu)^{\top}\Sigma^{-1}(x-\mu) = \sum_{i=1}^{m}\frac{(x_i-\mu_i)^2}{\sigma_i^2}.
\]
Подставляем в \eqref{eq:mvn} и разносим множители по компонентам, пользуясь тем,
что экспонента суммы есть произведение экспонент:
\[
  \dens(x)
  = \frac{1}{(2\pi)^{m/2}\prod_i\sigma_i}
    \exp\Bigl(-\frac12\sum_{i=1}^{m}\frac{(x_i-\mu_i)^2}{\sigma_i^2}\Bigr)
  = \prod_{i=1}^{m}\frac{1}{\sigma_i\sqrt{2\pi}}
    \exp\Bigl(-\frac{(x_i-\mu_i)^2}{2\sigma_i^2}\Bigr),
\]
то есть произведение плотностей $\Normal(\mu_i,\sigma_i^2)$. $\blacksquare$

Верно и обратное: если плотность распадается в произведение одномерных, то
$\Sigma$ диагональна и компоненты независимы. Для нормального распределения
некоррелированность и независимость -- это одно и то же (в общем случае, как
показано в задании 7, нет).

\zad{7}
\dano Компоненты $\xi_i$ и $\xi_j$ случайного вектора независимы и имеют
конечные вторые моменты.

\treb Показать, что $\sigma_{ij} = 0$.

\resh Центрированные величины $\xi_i - \mu_i$ и $\xi_j - \mu_j$ тоже независимы:
каждая является функцией только от своей исходной величины. Применяя
\eqref{eq:indep},
\[
  \sigma_{ij} = \Exp\bigl((\xi_i-\mu_i)(\xi_j-\mu_j)\bigr)
              = \Exp(\xi_i-\mu_i)\cdot\Exp(\xi_j-\mu_j)
              = 0\cdot 0 = 0 . \qquad\blacksquare
\]
Условие конечности вторых моментов нужно, чтобы ковариация вообще была
определена.

\smallskip
\emph{Обратное неверно.} Пусть $\xi \sim \mathrm{U}[-1,1]$ и $\eta = \xi^2$.
Величины зависимы предельно жёстко: значение $\eta$ полностью определяется
значением $\xi$. Тем не менее
\[
  \mathrm{cov}(\xi,\eta) = \Exp(\xi\cdot\xi^2) - \Exp\xi\cdot\Exp\xi^2
  = \Exp\xi^3 = \int_{-1}^{1}\frac{x^3}{2}\dd x = 0,
\]
так как подынтегральная функция нечётна. Ковариация улавливает только линейную
связь. Отсюда важное следствие: метод главных компонент делает признаки
некоррелированными, но не независимыми.

\zad{8}
\dano Ковариационная матрица $\Sigma$ случайного вектора $\xi$ с вектором
средних $\mu$.

\treb Доказать, что $\Sigma$ неотрицательно определена, и выяснить, в каком
случае она не является положительно определённой.

\zachem Знак определённости $\Sigma$ решает, существует ли многомерная нормальная
плотность \eqref{eq:mvn}: при $\det\Sigma = 0$ обратной матрицы нет.

\resh Напомним определения: матрица $A$ неотрицательно определена, если
$a^{\top}Aa \ge 0$ для всех $a$, и положительно определена, если неравенство
строгое при $a \ne 0$.

Возьмём произвольный вектор $a \in \R^m$ и построим по нему скалярную величину
$\eta = a^{\top}(\xi - \mu) = \sum_i a_i(\xi_i - \mu_i)$. Раскроем квадратичную
форму, подставив определение $\sigma_{ij}$:
\[
  a^{\top}\Sigma a
  = \sum_{i,j} a_i a_j \Exp\bigl((\xi_i-\mu_i)(\xi_j-\mu_j)\bigr)
  = \Exp\Bigl(\sum_i a_i(\xi_i-\mu_i)\Bigr)^{\!2}
  = \Exp\eta^2 \ge 0 .
\]
Внести сумму под знак математического ожидания можно по линейности, сумма
конечна. Симметричность очевидна из $\sigma_{ij} = \sigma_{ji}$. $\blacksquare$

Теперь второй вопрос. Положительной определённости нет ровно тогда, когда для
некоторого $a \ne 0$ получается $\Exp\eta^2 = 0$. Неотрицательная величина имеет
нулевое среднее только если она сама равна нулю почти наверное, поэтому
\[
  \sum_{i=1}^{m} a_i(\xi_i - \mu_i) = 0 \ \text{п.\,н.}
  \iff a^{\top}\xi = a^{\top}\mu = \mathrm{const} \ \text{п.\,н.}
\]
Это означает, что компоненты вектора признаков линейно зависимы: всё
распределение сосредоточено на гиперплоскости размерности меньше $m$. Тогда
$\det\Sigma = 0$, матрицы $\Sigma^{-1}$ не существует и плотность
\eqref{eq:mvn} не определена.

На практике так бывает, когда признак продублирован, когда один признак является
суммой или другой линейной комбинацией остальных, или когда элементов выборки
меньше, чем признаков. То же условие делает вырожденной систему нормальных
уравнений в методе наименьших квадратов.

\zad{9}
\dano Полная группа попарно несовместных событий $W_1,\dots,W_n$ с
$\PP(W_i) > 0$ (возможные предпосылки) и событие $X$ с $\PP(X) > 0$.

\treb Доказать формулу Байеса и показать, что сумма апостериорных вероятностей
всех предпосылок равна единице.

\resh \emph{Формула Байеса.} Запишем вероятность пересечения $W_i \cap X$ двумя
способами через определение условной вероятности:
\[
  \PP(W_i \mid X)\,\PP(X) = \PP(W_i \cap X) = \PP(X \mid W_i)\,\PP(W_i).
\]
Делим на $\PP(X) > 0$:
\begin{equation}\label{eq:bayes}
  \PP(W_i \mid X) = \frac{\PP(X \mid W_i)\,\PP(W_i)}{\PP(X)} . \qquad\blacksquare
\end{equation}
Никаких дополнительных предположений не потребовалось: формула Байеса -- это
определение условной вероятности, записанное в другую сторону.

\emph{Сумма апостериорных вероятностей.} Подставим в \eqref{eq:bayes} формулу
полной вероятности \eqref{eq:total} и просуммируем по $i$:
\[
  \sum_{i=1}^{n}\PP(W_i \mid X)
  = \frac{\sum_{i=1}^{n}\PP(X \mid W_i)\PP(W_i)}{\PP(X)}
  = \frac{\PP(X)}{\PP(X)} = 1 . \qquad\blacksquare
\]
Числитель и знаменатель оказались одним и тем же выражением, потому что
знаменатель как раз и собран из всех числителей.

\smallskip
Отсюда практическое следствие для классификации. Числители
$\PP(X \mid W_i)\PP(W_i)$ сами по себе в сумму единицы не складываются, единицу
даёт деление на $\PP(X)$. Но при сравнении классов для одного и того же
наблюдения знаменатель у всех одинаков, поэтому решающее правило можно строить
по числителю, не вычисляя $\PP(X)$ вовсе.

\end{document}
